\section{Complete flag reference}\label{app:flags}

This appendix lists every flag accepted by \CJKQC{} 0.1.0. Flags are grouped
by purpose; within a group the order is not significant. Words in angle
brackets are placeholders you replace with your own value, and square brackets
mark optional values.

\subsection{Inputs and output}

\begin{longtable}{@{}L{4.0cm}L{9.5cm}@{}}
\toprule
\textbf{Flag} & \textbf{Meaning} \\
\midrule
\endhead
\texttt{INPUTS...} & Positional, required. Program outputs (build mode) or
record databases (merge mode). \\
\flg{-o}, \flg{--out <PATH>} & Database to write. Required unless a
manipulation or print flag makes it optional. \\
\flg{--fout} & Force replacement of an existing \flg{--out} or
\texttt{movie.xyz} without asking. \\
\flg{--head-bytes <N>} & How much of the start of each program output is read,
in bytes (default 262144, about 256 kilobytes). \\
\flg{--tail-bytes <N>} & How much of the end of each program output is read,
in bytes (default 4194304, about 4 megabytes). \\
\flg{--parallel <N>} & How many files are read at once (default 1). Merge
mode ignores it. \\
\flg{--movie} & Write final structures to \texttt{movie.xyz}, one frame per
\fld{ok} record with a geometry. \\
\bottomrule
\end{longtable}

\subsection{Print columns}

There is no default column selection. Columns always appear in the fixed
order of Section~\ref{sec:printing}.

\begin{longtable}{@{}L{3.5cm}L{10.0cm}@{}}
\toprule
\textbf{Flag} & \textbf{Column and quantity} \\
\midrule
\endhead
\flg{--filename} & \fld{filename}: the record's \fld{name} (file stem). \\
\flg{--name} & \fld{name}: composition written without separators
(\texttt{2am3sa}). \\
\flg{--path} & \fld{path}: the full path of the original output file. \\
\flg{--status} & \fld{status}: \fld{ok}, \fld{failed}, or \fld{truncated}. \\
\flg{--program} & \fld{program}: program name and version (\texttt{xtb 6.7.1}). \\
\flg{--command} & \fld{command}: the run's command line, double-quoted. \\
\flg{--correction} & \fld{correction}: correction kind and source details. \\
\flg{--el} & \fld{el}: electronic energy. \\
\flg{--zpe} & \fld{zpe}: harmonic zero-point energy. \\
\flg{--h} & \fld{h}: enthalpy at \flg{--temp}. \\
\flg{--s} & \fld{s}: entropy term $T\,S$ at \flg{--temp}. \\
\flg{--g} & \fld{g}: Gibbs free energy at \flg{--temp}. \\
\flg{--elcorr} & \fld{elcorr}: attached correction energy. \\
\flg{--hcorr} & \fld{hcorr}: enthalpy from the correction energy. \\
\flg{--gcorr} & \fld{gcorr}: Gibbs free energy from the correction energy. \\
\flg{--pop [LEVEL]} & Population column; \texttt{LEVEL} is \fld{g} (default),
\fld{el}, \fld{elcorr}, or \fld{gcorr}. Repeatable. \\
\flg{--rg} & \fld{rg}: mass-weighted radius of gyration in \AA{}. \\
\flg{--mass} & \fld{mass}: total mass in u. \\
\flg{--dip} & \fld{dip}: dipole magnitude in atomic units. \\
\flg{--dips} & \fld{dips}: dipole vector components in atomic units. \\
\flg{--pol} & \fld{pol}: isotropic polarizability in atomic units. \\
\flg{--pols} & \fld{pols}: polarizability tensor values in atomic units. \\
\flg{--charge} & \fld{charge}: net charge. \\
\flg{--multiplicity} & \fld{multiplicity}: spin multiplicity. \\
\flg{--rsn} & \fld{rsn}: rotational symmetry number. \\
\flg{--natoms} & \fld{natoms}: geometry atom count. \\
\flg{--lf} & \fld{lf}: lowest stored frequency in cm$^{-1}$. \\
\flg{--freqs} & \fld{freqs}: all stored frequencies in cm$^{-1}$. \\
\flg{--nfreq} & \fld{nfreq}: number of stored frequencies. \\
\flg{--nimag} & \fld{nimag}: number of negative frequencies. \\
\flg{--rots} & \fld{rots}: rotational constants in cm$^{-1}$. \\
\flg{--components} & \fld{components}: composition as \texttt{2am+3sa}. \\
\bottomrule
\end{longtable}

\subsection{Filtering and manipulation}

\begin{longtable}{@{}L{4.3cm}L{9.2cm}@{}}
\toprule
\textbf{Flag} & \textbf{Meaning} \\
\midrule
\endhead
\flg{--extract <QUERY>} & Keep records whose composition matches the query. \\
\flg{--nextract <QUERY>} & Drop records whose composition matches the query. \\
\flg{--rename <OLD> <NEW>} & Move the count of component code \texttt{OLD} to
code \texttt{NEW}. Repeatable; renames apply in the order given. \\
\flg{--cut <PROP> <V>} & Absolute cutoff, property $\le V$. \\
\flg{--filter-lt <PROP> <V>} & Absolute cutoff, property $< V$. \\
\flg{--pass <PROP> <V>} & Absolute cutoff, property $> V$. \\
\flg{--filter-ge <PROP> <V>} & Absolute cutoff, property $\ge V$. \\
\flg{--is <PROP> <V>} & Absolute cutoff, property $= V$. \\
\flg{--isnot <PROP> <V>} & Absolute cutoff, property $\ne V$. \\
\flg{--cutr <PROP> <V>} & Relative cutoff, property $\le V$ above the group
minimum. \\
\flg{--rel-filter-lt <PROP> <V>} & Relative cutoff, $<V$ above the minimum. \\
\flg{--passr <PROP> <V>} & Relative cutoff, $>V$ above the minimum. \\
\flg{--rel-filter-ge <PROP> <V>} & Relative cutoff, $\ge V$ above the
minimum. \\
\flg{--uniq <MODE>} & \texttt{dup}, or a comma list of
\fld{metric[decimals]} properties. \\
\flg{--reacted} & Keep the composition group's majority fragment pattern. \\
\flg{--reacted2} & Keep only the deviating fragment patterns. \\
\flg{--sort <PROPERTY>} & Sort by \fld{el}, \fld{g}, \fld{elcorr}, \fld{gcorr},
\fld{rg}, \fld{mass}, or \fld{filename}. \\
\flg{--reverse} & Reverse the sort orientation. \\
\flg{--select <N>} & Keep the first $N$ records of each group after sorting. \\
\flg{--glob} & Keep each group's lowest-Gibbs record. \\
\flg{--shuffle} & Put the final result in a random order; the same seed gives
the same order. \\
\flg{--seed <N>} & Seed for \flg{--shuffle} (default 0). \\
\flg{--noname} & Treat the whole database as one group; \flg{--extract}
compares names exactly. \\
\bottomrule
\end{longtable}

\begin{note}
The help text for \flg{--sort} mentions ``name'', but the accepted value is
\fld{filename}, which sorts by the record's \fld{name}. Use
\flg{--sort filename}.
\end{note}

\subsection{Thermodynamics}

\begin{longtable}{@{}L{3.5cm}L{10.0cm}@{}}
\toprule
\textbf{Flag} & \textbf{Meaning} \\
\midrule
\endhead
\flg{--temp <K>} & Temperature in K (default 298.15; must be $>0$). \\
\flg{--fc <CM-1>} & Quasi-harmonic reference frequency (default 100; 0 selects
pure harmonic entropy). \\
\flg{--fscale <F>} & Vibrational frequency scale factor (default 1.0; must be
$>0$). \\
\flg{--makereal} & Treat imaginary frequencies as their absolute values
inside the derivation. \\
\flg{--ignoreimag} & Drop imaginary frequencies from the vibrational sums.
Conflicts with \flg{--makereal}. \\
\flg{--unit} & Print energies in kcal/mol instead of hartree. Takes no value. \\
\bottomrule
\end{longtable}

\subsection{Corrections and extrapolation}

\begin{longtable}{@{}L{4.3cm}L{9.2cm}@{}}
\toprule
\textbf{Flag} & \textbf{Meaning} \\
\midrule
\endhead
\flg{--add-correction <PATH>} & Attach single-point corrections from a
database, matching by exact case-sensitive \fld{name}. One path; not
repeatable. \\
\flg{--extrapolate <MODE VALUES>...} & Attach an extrapolated correction.
Modes: \texttt{2p LOW HIGH N ALPHA BETA}, \texttt{2s LOW HIGH F}, or
\texttt{2p LOWPAIR HIGHPAIR N ALPHA BETA F}. Place it after the positional
inputs. \\
\bottomrule
\end{longtable}

\subsection{Formation}

\begin{longtable}{@{}L{4.3cm}L{9.2cm}@{}}
\toprule
\textbf{Flag} & \textbf{Meaning} \\
\midrule
\endhead
\flg{--formation} & Print the raw formation-energy section. \\
\flg{--formation-corr} & Print the corrected formation-energy section. \\
\flg{--conc <CODE> <VALUE>} & Evaluate formation free energies at a vapor
concentration. Repeatable; one setting per code. \\
\flg{--cnt} & Use the self-consistent $(N-1)/N$ concentration scaling. \\
\flg{--bavg} & Collapse each printed formation section into one
Boltzmann-averaged row per composition. \\
\bottomrule
\end{longtable}

\subsection{Solvation and hydration}

\begin{longtable}{@{}L{4.3cm}L{9.2cm}@{}}
\toprule
\textbf{Flag} & \textbf{Meaning} \\
\midrule
\endhead
\flg{--hydration} & Print the raw water distribution section (solvent
\texttt{w}). \\
\flg{--hydration-corr} & Print the corrected water distribution section. \\
\flg{--solvation <CODE>} & Print the raw distribution section for solvent
\texttt{CODE}. \\
\flg{--solvation-corr <CODE>} & Print the corrected distribution section for
solvent \texttt{CODE}. \\
\flg{--rh <VALUE>} & Relative humidity as a fraction (0.5 is 50 \%). Water
only. \\
\flg{--psolvent <VALUE>} & Solvent partial pressure in the \flg{--conc}
VALUE grammar. Overrides \flg{--rh}; required for non-water solvents. \\
\bottomrule
\end{longtable}

\subsection{Help and version}

\begin{longtable}{@{}L{3.5cm}L{10.0cm}@{}}
\toprule
\textbf{Flag} & \textbf{Meaning} \\
\midrule
\endhead
\flg{-h}, \flg{--help} & Print help and exit. \\
\flg{-V}, \flg{--version} & Print the version and exit. \\
\bottomrule
\end{longtable}
